\section{Discussion}
\label{sec:discussion}

\subsection{Why smoothing helps and the dynamic program does not}
\label{sec:disc-why}

The two temporal results are not in tension; they follow from the same fact.
Split CP calibrates a \emph{scalar} per epoch and produces a set by
thresholding it. Any transformation of the model output that makes the true
label's scalar smaller relative to the other labels --- which the
forward--backward posterior does, because it integrates evidence over the
whole recording --- improves efficiency without touching validity. The
dynamic program does something different: it minimises over \emph{paths}, and
a path-level minimum is always at most a point-level one, so the resulting
level comparison is looser rather than tighter. The transition structure that
the DP is meant to exploit has already been absorbed into the smooth posterior
at the point where the score is computed. This is why the ablation at $\beta=0$
is the decisive one: it shows the two constructions are identical there, and
positive $\beta$ only adds looseness.

The practical consequence is a simpler recipe than the one we set out to
build: smooth, score the deficit, calibrate on subjects, adapt per patient.
That recipe needs no dynamic program and no block-level calibration.
It also separates two sources of ambiguity. Irreducible disagreement over an
N1 or transition epoch remains visible as a larger set, whereas a mismatch
between a new patient and the calibration cohort is corrected by the adaptive
stream. This is the operational counterpart of the aleatoric--epistemic
distinction discussed for sleep staging \citep{vangorp2022certainty}.

\subsection{What per-patient adaptation costs and buys}
\label{sec:disc-aci}

Per-patient ACI requires labels on a delay. In sleep staging the delay is the
technologist's review, so the feedback assumption is realistic, but it does
mean that the guarantee is about the \emph{realised} sequence of predictions
and not about a prediction made before any review has happened. The first
epochs of a new patient are therefore predicted with the pooled threshold,
and adaptation only begins to pay after a few dozen epochs. The coverage
vectors in \Cref{sec:results-aci} show the mechanism: easy patients lose
coverage that they did not need, hard patients gain coverage that they did,
and the aggregate changes by less than a point. This reallocation is the
point, not a side effect.

The step scale is the one hyperparameter, and it is chosen without test
labels by simulating held-out streams inside the calibration cohort. That
choice is conservative: it optimises the mean one-sided shortfall rather than
the worst patient, so a deployment that cares about the single worst patient
should select on that criterion instead; the released code exposes the choice
as \texttt{--criterion worst}.

\subsection{Reading the sets in a review workflow}
\label{sec:disc-reading}

A single coverage number is not enough to evaluate a set-valued stager, for
three reasons that this study makes concrete. It is marginal over epochs and
therefore blind to the patient whose sets are systematically too small. It is
satisfied trivially by large sets, so it must be read against set size and
review workload. And it says nothing about temporal behaviour, which in a
hypnogram is the difference between a usable display and a flickering one. We
would therefore argue that a conformal sleep-staging paper should report
per-patient coverage, the fraction of patients below nominal, set size, and a
temporal consistency statistic, and should use subject-wise folds that
partition the cohort.

Those four numbers translate directly into the review step. A night of eight
hours contains 960 epochs; at the $0.76$ workload of pooled APS, 730 of them
would be sent back for adjudication, whereas the smoothed deficit score sends
about 150, plus the empty sets that also return to the reviewer. The size of a
non-singleton set is itself informative, because it separates ``N2 or N3''
from ``any of four stages'': only the first is a question a technologist can
settle by re-reading a page of the record. It is worth being explicit that a
set of two stages is not automatically a weaker statement than a point label.
Inter-scorer studies show that two trained technologists frequently disagree
on N1 and on transition epochs \citep{rosenberg2013}, so an admitted ambiguity
of that size is commensurate with the reference standard against which any
stager is judged, whereas a confident point prediction that is wrong is not.

The resulting deployment pattern is short. The stager is frozen; its HMM
transitions come from the training labels; one pooled threshold is calibrated
on the handful of scored nights a unit already has; and the per-patient
adaptive stream runs during review, consuming each stage label after it has
been confirmed. Nothing in the procedure is fitted to the patient being
scored, which is the property that makes it usable in a unit that cannot
collect labelled epochs in advance.

\subsection{Limitations}
\label{sec:limitations}

\paragraph{Single channel, five classes}
All experiments use Fpz--Cz only and the five AASM stages. Multi-channel
recordings, and problems that use the finer R\&K stages, will have different
score distributions; the methods apply unchanged, but the specific numbers
should not be extrapolated.

\paragraph{Cohort size}
The main protocol tests four subjects per fold, and the merged estimate is
over 20 subjects. This is small, and the seed spread we report is a
spread over training runs rather than over patient cohorts. The extended
cohort in \Cref{sec:results-cohort} is a partial answer, but a
multi-centre study with several hundred patients would be needed to put a
confidence interval on the worst-patient statistic.

\paragraph{Feedback assumption}
ACI needs the label after the prediction. In a fully automatic deployment
with no review, the stream is unavailable and only the pooled or warm-up
rules apply.

\paragraph{N1}
The backbones used here reproduce the known weakness of the literature: N1
F1 is low. The consequence for the sets is asymmetric. Pooled APS
over-covers N1 because its threshold is loose, whereas the smoothed deficit
score omits N1 in more than half of the N1 epochs (coverage $0.43$ against
nominal $0.90$), because the model is confidently wrong on exactly those
epochs. This qualifies the efficiency claim of \Cref{sec:results-score};
thresholding the same score per stage restores every stage to nominal
coverage, and \Cref{tab:perclass} reports what that costs in set size
($1.74$ against $1.53$ labels). A stronger backbone would shift every number;
we expect, but have not shown, that it would leave the ordering of the methods
unchanged, because the ordering is driven by the protocol rather than by the
model.

\paragraph{No clinical evaluation}
We make no claim that a technologist working with these sets is faster or
more accurate than one working with point predictions. Review workload is a
proxy, and it should be validated in a reader study before deployment.

\section{Conclusions}
\label{sec:conclusion}

Conformal prediction is a natural way to give a clinician a statement about a
sleep-stage prediction that is honest under subject shift. We have shown
that the statement the standard method makes is the wrong one --- pooled over
epochs rather than per patient --- and that two simple changes together
address it.
Scoring
the smoothed posterior deficit shrinks the average set by about forty per cent
at comparable coverage and cuts the fraction of epochs needing manual
review from $0.76$ to $0.16$ (to $0.49$ if the threshold is conditioned on the
stage, which holds every stage at nominal), and running an independent
adaptive stream per patient moves coverage from easy patients to hard ones
without using any test labels for tuning. The dynamic-programming sequence constructors that
motivated this work turn out not to help, for a reason we can prove, and we
report that result explicitly, together with the exact protocol and code
needed to check it.

The practical form of the result is a stager that reports, epoch by epoch, the
stages it cannot rule out, and a number of epochs it cannot settle. That does
not replace the technologist's review; it tells the technologist where the
review is worth spending.
